\documentclass[11pt]{article}
\usepackage{amsmath,amssymb,amsthm}
\usepackage[margin=1in]{geometry}
\title{Problem 449: Chocolate Covered Candy}
\author{Project Euler Solutions}
\date{}

\begin{document}
\maketitle

\section{Problem Statement}
A candy has the shape of a body of revolution formed by rotating an ellipse about its minor axis, then coated with a uniform chocolate layer of thickness $t$. Find the volume of chocolate for an ellipse with semi-major axis $a$ and semi-minor axis $b$, to 8 decimal places.

\section{Mathematical Analysis}
The candy is an oblate spheroid (ellipsoid of revolution). The volume of an oblate spheroid with semi-major axis $a$ and semi-minor axis $b$ ($a > b$) is:
$$V_{\text{candy}} = \frac{4}{3}\pi a^2 b$$

The chocolate-covered candy is the parallel body (offset surface) at distance $t$ from the spheroid. The volume of a parallel body of a convex body $K$ at distance $t$ is given by Steiner's formula:
$$V(K_t) = V(K) + S(K) \cdot t + \pi M(K) \cdot t^2 + \frac{4}{3}\pi t^3$$

\section{Derivation}
For an oblate spheroid with semi-axes $a, a, b$ (with $a > b$):

\textbf{Volume:}
$$V = \frac{4}{3}\pi a^2 b$$

\textbf{Surface area:}
$$S = 2\pi a^2 + \frac{\pi b^2}{e} \ln\frac{1+e}{1-e}$$
where $e = \sqrt{1 - b^2/a^2}$ is the eccentricity.

\textbf{Mean width} $M$: For the Steiner formula, we need the mean curvature integral, which for an oblate spheroid is:
$$M = 2a + \frac{2b^2}{a \cdot e} \arcsin(e)$$

The chocolate volume is:
$$V_{\text{choc}} = V(K_t) - V(K) = S \cdot t + \pi M \cdot t^2 + \frac{4}{3}\pi t^3$$

For $a = 3, b = 1, t = 1$: answer $= 103.16172816$.

\section{Proof of Correctness}
Steiner's formula for parallel bodies of convex sets is a classical result in integral geometry. For smooth convex bodies, it follows from the tube formula and the principal curvatures of the boundary surface.

\section{Complexity Analysis}
\begin{itemize}
    \item \textbf{Analytical:} $O(1)$ --- closed-form computation.
    \item \textbf{Numerical:} High precision requires careful evaluation of $\ln$ and $\arcsin$ functions.
\end{itemize}

\section{Answer}

\[
\boxed{103.37870096}
\]

\end{document}
